\documentclass[11pt,letterpaper]{article}
\usepackage{fontspec}
\usepackage{tocloft}
\usepackage{multicol}
\usepackage{nicefrac}
\usepackage{url}
\usepackage[left=1.3in,right=1.3in,top=1in,bottom=1in]{geometry}
\setmainfont[Scale=1.1,AutoFakeBold=1.5,AutoFakeSlant=0.3]{Doves Type}

% Prevent page breaks within list items
\widowpenalty=10000
\clubpenalty=10000
\interlinepenalty=500

\title{Crispy-Skin Salmon with Lemon-Caper Brown Butter and Dijon Spinach}
\author{}
\date{}

\begin{document}

\maketitle
\thispagestyle{empty}

\begin{quote}\textit{Skin-on salmon fillets are seared hard in a hot pan and never flipped---the top cooks by spooning the pan's hot oil over it, so the skin stays in uninterrupted contact with the metal and crisps while the flesh sets evenly. Unsalted butter browned to a nutty \textit{beurre noisette}, then arrested with capers, lemon juice, and a splash of wine, becomes the sauce; shallots, garlic, and Dijon mustard wilt a pan of spinach alongside. The plate is a pool of spinach, crisp-skinned salmon, and the sauce spooned around it---never over the skin. Serves 2.}\end{quote}

\section*{Ingredients}
\setlength{\columnsep}{20pt}
\begin{multicols}{2}
\noindent
    Salmon fillets, (6~oz.), skin-on \dotfill 2 \\
    Kosher salt \dotfill 1\nicefrac{1}{2}~tsp. \\
    Olive oil \dotfill 2~Tbsp. \\
    Unsalted butter \dotfill 10~Tbsp. \\
    \columnbreak
    Capers, drained \dotfill 2~Tbsp. \\
    Lemon juice \dotfill 3~Tbsp. \\
    Dry white wine \dotfill \nicefrac{1}{4}~cup plus 2~Tbsp. \\
    Shallots \dotfill 1~large \\
    Garlic cloves \dotfill 2 \\
    Baby spinach, stemmed \dotfill 8~oz. \\
    Dijon mustard \dotfill 1~Tbsp. \\
    Fresh parsley, roughly chopped \dotfill 3~Tbsp. \\
    Black pepper, cracked \dotfill to taste \\
\end{multicols}

\section*{Directions}

\noindent
Blot \textbf{salmon} dry with paper towels; keep skin-side up ---
Finely chop \textbf{shallots} and mince \textbf{garlic}; combine in \textit{Small Bowl~\#1} (spinach aromatics) ---
Roughly chop \textbf{parsley}; set aside in \textit{Small Bowl~\#2} (sauce finish) ---
Drain \textbf{capers} into \textit{Small Bowl~\#3} ---
Combine 2~Tbsp. \textbf{white wine} with 2~Tbsp. \textbf{lemon juice} in \textit{Small Bowl~\#4} (sauce liquids) ---
Measure \nicefrac{1}{4}~cup \textbf{white wine} and 1~Tbsp. \textbf{lemon juice} for the \textbf{spinach}; have ready

\begin{enumerate}
    \item Blot the \textbf{salmon} skin dry once more and season both sides with \textbf{kosher salt} (about 1\nicefrac{1}{2}~tsp. total), sprinkling from a height so it falls in an even coat. Be generous---a bold seasoning now is what carries the fish. Set the fillets skin-side up on the board.

    \item Heat a 12-inch cast iron skillet over medium-high for \textit{3--4~minutes}, until you can hold your hand a few inches above the surface for only a moment. Add \textbf{olive oil} to coat the base of the pan and let it heat until it shimmers and moves loosely, about \textit{30~seconds}, then lower the heat to medium. A cool pan is the one thing that ruins this dish.

    \item Lay the \textbf{salmon} skin-side down away from you and press the fillets flat with your fingers for \textit{15--20~seconds} so the skin meets the metal across its whole surface; you will feel bubbles tickling your fingertips from inside the fish. It should sizzle hard on contact. Cook without moving for \textit{5~minutes}, until the flesh has turned opaque about a third of the way up the sides and a lifted corner of skin shows deep gold.

    \item Baste the fillets without flipping them: tilt the pan toward you and spoon the pooled hot oil continuously over the top for \textit{3--4~minutes}, until the flesh is opaque throughout with no translucent pink at the thickest point (the loin side) and the centre reaches \textit{140--145°F} on an instant-read thermometer. The skin stays in contact with the pan the entire time and finishes an even, crusty deep brown---that is the point of not flipping. Transfer the \textbf{salmon} to a plate, skin-side up, and let it rest uncovered while you make the sauce and the spinach---about \textit{8~minutes} in all, which finishes the carryover and lets the muscle relax. Leave it uncovered; a foil tent would steam the skin you just spent the sear on.

    \item Make the brown butter sauce in the same skillet. Pour off all but about 1~Tbsp. of the oil so the sauce tastes of butter rather than oil.
\begin{enumerate}
    \item Add 8~Tbsp. \textbf{butter} and melt over medium heat, swirling. Cook \textit{2--4~minutes} until the foam subsides and the milk solids at the bottom turn nut brown and smell of hazelnuts. In a black-bottomed pan the colour is hard to read; lift a spoonful into the light to judge it.
    \item Add \textbf{capers} (\textit{Small Bowl~\#3}) and stir; the moisture in them fizzles and stops the butter browning further.
    \item Add \textbf{white wine} and \textbf{lemon juice} (\textit{Small Bowl~\#4}). Swirl and scrape the pan for \textit{30~seconds}; the liquid cools the fat and the sauce comes together glossy and sharp. Slide the pan off the heat and keep it warm.
\end{enumerate}

    \item Cook the Dijon spinach in a second large skillet.
\begin{enumerate}
    \item Melt the remaining 2~Tbsp. \textbf{butter} over medium heat. Add \textbf{shallots} and \textbf{garlic} (\textit{Small Bowl~\#1}) with a pinch of \textbf{kosher salt} straight away---the salt pulls their moisture and starts them cooking. Sweat for \textit{2--3~minutes} until fragrant and softened but not browned.
    \item Add \textbf{baby spinach} with another small pinch of \textbf{kosher salt}. Toss with the aromatics and let it wilt in contact with the pan for \textit{2--3~minutes}; a lid speeds this but is not needed. Do not over-wilt---the leaves should be tender and dark green, not slumped.
    \item Add \nicefrac{1}{4}~cup \textbf{white wine} and 1~Tbsp. \textbf{Dijon mustard} and stir through for \textit{1~minute}. Finish off the heat with 1~Tbsp. \textbf{lemon juice} and \textbf{black pepper} to taste. The spinach is done when it stops shrinking and holds its shape.

\end{enumerate}

    \item Finish and plate. Stir \textbf{parsley} (\textit{Small Bowl~\#2}) into the warm---no longer hot---brown butter sauce and taste for salt; the parsley goes in last so it stays bright green. Spoon a pool of \textbf{spinach} onto each plate, lay a rested fillet beside it skin-side up, and spoon the caper brown butter around the plate and over the spinach, never over the skin. Draw the sauce and spinach together at the table and serve immediately.
\end{enumerate}

\newpage

{\footnotesize
\setlength{\columnsep}{20pt}
\setlength{\multicolsep}{6pt}
\begin{multicols}{2}
\setlength{\parindent}{0pt}
\setlength{\parskip}{2pt}
\setlength{\itemsep}{0pt}
\setlength{\parsep}{0pt}

\subsection*{Yield}
\begin{itemize}
    \item Serves 2 (one 6~oz. fillet each)
\end{itemize}

\subsection*{Equipment Required}
\begin{itemize}
    \item 12-inch cast iron or other heavy-bottomed skillet
    \item Second large skillet for the \textbf{spinach}
    \item Large spoon for basting
    \item Thin, flexible spatula
    \item Instant-read thermometer
    \item Long, sharp chef's knife and cutting board
    \item Plate or rack for resting the \textbf{salmon}
    \item Paper towels
    \item Small prep bowls (4)
    \item Measuring cups and spoons
\end{itemize}

\subsection*{Mise en Place}
\begin{itemize}
    \item Small Bowl \#1 --- \textbf{shallots} and \textbf{garlic}, chopped (spinach aromatics)
    \item Small Bowl \#2 --- \textbf{parsley}, roughly chopped (sauce finish)
    \item Small Bowl \#3 --- \textbf{capers}, drained
    \item Small Bowl \#4 --- 2~Tbsp. \textbf{white wine} combined with 2~Tbsp. \textbf{lemon juice} (sauce liquids)
    \item Measure \nicefrac{1}{4}~cup \textbf{white wine} and 1~Tbsp. \textbf{lemon juice} for the \textbf{spinach}; have ready
    \item Divide the \textbf{butter}: 8~Tbsp. for the sauce, 2~Tbsp. for the \textbf{spinach}
    \item Stem the \textbf{spinach} and keep the \textbf{salmon} skin-side up; blot the skin dry just before it goes into the pan
\end{itemize}

\subsection*{Ingredient Tips}
\begin{itemize}
    \item \textbf{Salmon}: Skin-on and fresh matter more than species---look for a bright, shiny surface, a clean sea smell, and flesh that pushes back when pressed. Dull, gray, spongy fillets will not crisp
    \item \textbf{Butter}: Unsalted and of good quality; the butter is the sauce, so its flavour is the whole sauce
    \item \textbf{Capers}: Small nonpareil capers are the right size; rinse them for less salt and chop any oversized ones
    \item \textbf{White wine}: Dry and unoaked---Muscadet, Sauvignon Blanc, or Chablis. Avoid anything sweet
    \item \textbf{Olive oil}: A lightly flavoured or unflavoured oil works too; the searing oil is poured off before the sauce
    \item \textbf{Spinach}: Baby spinach needs only stemming; mature spinach works but stem it fully and allow \textit{1--2~minutes} more wilt time
    \item \textbf{Garlic}: Fresh whole cloves only---pre-peeled garlic is a different, flatter ingredient
    \item \textbf{Parsley}: Flat-leaf, chopped just enough to tame it; a little stem is fine
\end{itemize}

\subsection*{Preparation Tips}
\begin{itemize}
    \item Moisture is the enemy of crisp skin: store the \textbf{salmon} skin-side up in the fridge so air circulates over it, and blot it dry twice---once when portioning, once before seasoning
    \item Season from a height; the salt spreads in an even coat instead of landing in patches
    \item Never crowd the pan; two fillets in a 12-inch skillet is the maximum, and they should sizzle loudly on contact
    \item Press the fillets down at the start and do not flip them; the skin must stay on the metal the whole time
    \item Judge the brown butter by smell and colour rather than time, and be ready to arrest it with the \textbf{capers} and the \textbf{lemon}-and-\textbf{wine} splash the moment it turns nut brown
    \item Salt the \textbf{shallots} and \textbf{garlic} the moment they hit the butter; it draws their moisture and removes the raw edge without browning them
    \item Pull the \textbf{salmon} at \textit{140--145°F} and expect carryover of a few degrees while it rests. The USDA recommends \textit{145°F} for salmon; at this point the flesh is opaque throughout and still moist
\end{itemize}

\subsection*{Make Ahead \& Storage}
\begin{itemize}
    \item Salt the \textbf{salmon} and hold it skin-side up, uncovered, in the fridge for up to \textit{1~hour} before cooking; the surface dries and crisps faster
    \item Chop the \textbf{shallots}, \textbf{garlic}, and \textbf{parsley} and stem the \textbf{spinach} up to \textit{4~hours} ahead; refrigerate
    \item The sauce and the \textbf{spinach} are best eaten immediately; the sauce separates and the spinach slumps on reheating
    \item Leftovers keep \textit{1--2~days} refrigerated; reheat gently---the skin will not stay crisp
\end{itemize}

\subsection*{Serving Suggestions}
\begin{itemize}
    \item A crusty baguette at the table to soak up the brown butter and lemon
    \item A cold, dry white---Chablis or Muscadet---or a light red served cool
    \item Roasted or smashed potatoes make a plain, absorbent partner; rice works too
\end{itemize}

\subsection*{Source}
\begin{itemize}
    \item Adapted from Ed Szymanski (Dame and Lord's, NYC), via Bon App\'etit \textit{Made in House}: \url{https://www.youtube.com/watch?v=nBvoe1AjNAk}
\end{itemize}

\end{multicols}
}

\end{document}
